% NichidaiMasterExam_R8_1st_1
\begin{tcolorbox} %R8_1st_1
    $\fbox{1}$
\begin{equation*}
    A=
    \begin{pmatrix}
        4& -1& 2\\
        -6& 3& -6\\
        -6& 2& -4
    \end{pmatrix}
\end{equation*}
    とする。また実線型写像$f\colon \R^3\to \R^3$を$f(x)= Ax\ (x\in \R^3)$で定める。次の問いに答えよ。
\begin{enumerate}
    \item $A$の行列式を求めよ。
    \item $f$の核と像の基底をそれぞれ求めよ。また$f$の核と次元をそれぞれ求めよ。
    \item $A$の固有値を全て求めよ。
    \item $\R^3$の基底$E$で、$f$の$E$に関する表現行列が$\begin{pmatrix}\lambda_1& 0& 0\\ 0& \lambda_2& 0\\ 0& 0& \lambda_3\end{pmatrix}$となるようなものを1つ求めよ。
    \item $\R^3$の基底$F$で、$f$の$F$に関する表現行列の成分がすべて偶数となるようなものはないことを示せ。
\end{enumerate}
\end{tcolorbox} 

\begin{souanswer} %R8_1st_1_ans
\begin{enumerate}
  \item 第一列について基本変形を行う．
\begin{equation*}
  \det{A}= 
\begin{vmatrix}
  4& -1& 2//
  -6& 3& -6\\
  -6& 2& -4
\end{vmatrix}= 
\begin{vmatrix}
  4& -1& 2\\
  0& 1& -2\\
  -2& 1& -2
\end{vmatrix}= 0
\end{equation*}
  \item $A$を行基本変形して階段行列をつくる．
\begin{equation*}
  A= 
\begin{vmatrix}
  4& -1& 2//
  -6& 3& -6\\
  -6& 2& -4
\end{vmatrix}= 
\begin{vmatrix}
  1& 0& 0//
  0& 1& -2\\
  0& 0& 0
\end{vmatrix}
\end{equation*}
  よって階数は$\rank{A}= 2$．$A\bm{x}= \bm{0}$を解くと，
\begin{equation*}
\begin{pmatrix}
  1& 0& 0\\
  0& 1& -2\\
  0& 0& 0
\end{pmatrix}\begin{pmatrix}x_1\\ x_2\\ x_3\end{pmatrix}= \begin{pmatrix}0\\ 0\\ 0\end{pmatrix}
  \iff \begin{cases}x_1= 0\\ x_2= 2x_3\end{cases}
\end{equation*}
  $x_3= t$とおくと，$\begin{pmatrix}x_1\\ x_2\\ x_3\end{pmatrix}= t\begin{pmatrix}0\\ 2\\ 1\end{pmatrix}$．\\
\begin{itemize}
  \item $\ker{f}$の基底：$\begin{pmatrix}0\\ 2\\ 1\end{pmatrix}$
  \item $\ker{f}$の次元：$\dim{\ker{f}}= 1$
\end{itemize}
  階段行列の主成分は第一列と第二列にあるため，$A$の第一列と第二列のベクトルが$\img{f}$の基底を作る．
\begin{itemize}
  \item $\img{f}$の基底：$\begin{Bmatrix}\begin{pmatrix}4\\ -6\\ -6\end{pmatrix}, \begin{pmatrix}-1\\ 3\\ 2\end{pmatrix}\end{Bmatrix}$
  \item $\img{f}$の次元：$\dim{\img{f}}= \rank{A}= 2$
\end{itemize}
  \item $A$の固有多項式$\det{\lambda I- A}= 0$を計算する．
\begin{equation*}
  \det{\lambda I- A}= 
\begin{vmatrix}
  \lambda- 4& 1& -2\\
  6& \lambda- 3& 6\\
  6& -2& \lambda+ 4
\end{vmatrix}= \lambda(\lambda- 1)(\lambda- 2)
\end{equation*}
  $\det{\lambda I- A}= 0$を解くと，$\lambda= 0, 1, 2$
  \item 各固有値に対して固有ベクトルを求める．
\begin{enumerate_roman}
  \item $\lambda= 0$のとき：$\bm{v}_{\lambda= 0}= \begin{pmatrix}0\\ 2\\ 1\end{pmatrix}$
  \item $\lambda= 1$のとき：$(I- A)\bm{v}_{\lambda= 1}= \bm{0}$を解く．
\begin{equation*}
  (I- A)\bm{v}_{\lambda= 1}= \begin{pmatrix}-3& 1& -2\\ 6& -2& 6\\ 6& -2& 5\end{pmatrix}\bm{v}_{\lambda= 1}= \bm{0}\iff \bm{v}_{\lambda= 1}= \begin{pmatrix}1\\ 3\\ 0\end{pmatrix}
\end{equation*}
  \item $\lambda= 2$のとき：$(2I- A)\bm{v}_{\lambda= 2}= \bm{0}$を解く．
\begin{equation*}
  (2I- A)\bm{v}_{\lambda= 2}= \begin{pmatrix}-2& 1& -2\\ 6& -1& 6\\ 6& -2& 6\end{pmatrix}\bm{v}_{\lambda= 2}= \bm{0}\iff \bm{v}_{\lambda= 2}= \begin{pmatrix}-1\\ 0\\ 1\end{pmatrix}
\end{equation*}
  求める基底$E$の1つは$\begin{Bmatrix}\begin{pmatrix}0\\ 2\\ 1\end{pmatrix}, \beign{pmatrix}1\\ 3\\ 0\end{pmatrix}, \begin{pmatrix}-1\\ 0\\ 1\end{pmatrix}\end{Bmatrix}$．このとき表現行列は$\begin{pmatrix}0& 0& 0\\ 0& 1& 0\\ 0& 0& 2\end{pmatrix}$となる．
\end{enumerate_roman}
  \item 背理法で示す．ある基底$F$に関する$f$の表現行列を$B= (b_{ij})\in M_3(\R)$とし，$B$のすべての成分$b_{ij}\ (1\le i, j\le 3)$が偶数整数であると仮定する．標準基底から基底$F$の基底変換行列を正則行列$P$とおくと，表現行列の変換規則より
\begin{equation*}
  B= P^{-1}AP
\end{equation*}
  が成り立つ．すなわち$A$と$B$は相似．相似な正方行列のトレースは不変であるため，
\begin{equation*}
  \tr{A}= \tr{B}
\end{equation*}
  が成り立つ．
  ここで仮定より$B$のすべての対角成分$b_{11}, b_{22}, b_{33}$は偶数である．したがってその和も偶数となる．\\
  一方行列$A$のトレースは
\begin{equation*}
  \tr{A}= 4+ 3+ (-4)= 3
\end{equation*}
  となるため，矛盾．\\
  したがってすべての成分が偶数となるような表現行列をもつ基底$F$は存在しない．
\end{enumerate}
\end{souanswer}
